\chapter{Estado del arte y fundamentos}
\label{chap:estado}

Este capítulo resume el contexto técnico del trabajo: primero los mecanismos de transmisión de
radio variable, que son el objeto de estudio, y después las tecnologías con las que se ha
construido el sistema de ensayo (accionamiento, medida de fuerza, microcontrolador, buses de
comunicación y supervisión).

\section{Mecanismos de transmisión con radio variable}

\subsection{Del \emph{fusee} a los tambores de radio variable}

El problema de convertir una fuerza variable en un par constante es antiguo. En los relojes de
muelle, el par del muelle real disminuye a medida que se desenrolla, y eso altera la marcha. El
\emph{fusee} lo resolvió con un cono de perfil curvo y una ranura helicoidal por la que se enrolla
una cadena: cuando el muelle está tenso, la cadena tira del radio pequeño del cono, y cuando se
debilita, del radio grande, de manera que el par entregado al tren de engranajes se mantiene
aproximadamente constante \cite{landes1983}. La condición de diseño es simplemente que el producto
fuerza por radio se conserve a lo largo del recorrido.

La misma idea se ha formalizado en la robótica bajo el nombre de tambores de radio variable
(\emph{Variable Radius Drums}, VRD). Seriani y Gallina \cite{seriani2016} plantean la síntesis del
perfil de un tambor para que una carga siga una ley de fuerza prescrita, y obtienen soluciones
analíticas cerradas para varios casos. En los robots paralelos accionados por cables, Bieber
\emph{et al.} \cite{bieber2023} diseñan el perfil del tambor a partir del perfil de carga máxima en
el espacio de trabajo, de manera que el par del motor no supere un valor dado para ninguna longitud
de cable. La fabricación aditiva ha abaratado mucho la producción de estos perfiles, lo que explica
el interés reciente.

\subsection{Compensación de gravedad con poleas no circulares}

Un segundo campo de aplicación es la compensación estática de peso. Herder \cite{herder2001}
estudia de forma sistemática los mecanismos estáticamente equilibrados, en los que la energía
potencial total se mantiene constante y el actuador sólo tiene que vencer las pérdidas. Endo
\emph{et al.} \cite{endo2010} proponen una polea no circular combinada con un muelle para
compensar el par gravitatorio de un brazo, y deducen numéricamente la forma de la polea. Shen y
Yano \cite{shen2022} refinan el cálculo de una polea en espiral teniendo en cuenta la variación de
la tensión del cable, y mejoran así la calidad de la compensación.

El mecanismo de este trabajo comparte el principio con todos ellos, pero con una diferencia: aquí
no hay un muelle que almacene energía, sino un motor que aporta todo el trabajo de elevación. La
caracola no reduce la energía necesaria para subir la masa; lo que hace es repartirla de forma
uniforme a lo largo del giro, para que el par de pico sea el menor posible. El
capítulo~\ref{chap:modelo} demuestra que ese reparto uniforme es precisamente el óptimo.

\section{Accionamiento: servomotores paso a paso de lazo cerrado}

Un motor paso a paso convencional funciona en lazo abierto: el controlador impone la secuencia de
fases y se da por hecho que el rotor la sigue. Es sencillo y tiene buen par a baja velocidad, pero
si la carga supera el par disponible se pierden pasos sin que el sistema lo detecte
\cite{kenjo1994}. Los servomotores paso a paso de lazo cerrado añaden un codificador al motor y un
accionamiento que regula la corriente en función del error de posición. Conservan el par a baja
velocidad del paso a paso, no pierden pasos y ofrecen varios modos de control (posición, velocidad
o par), además de magnitudes de supervisión como la velocidad real o el par entregado.

El accionamiento usado en este trabajo pertenece a esta familia
\completar{marca y modelo del motor y del accionamiento, par nominal, tensión de alimentación,
resolución del codificador}. Integra una interfaz RS-485 con protocolo Modbus RTU, a través de la
cual se configura (modo de control, rigidez, rampas, resistencia de frenado), se habilita, se le
envía la consigna de velocidad y se leen la velocidad y el par. Tener todo en un único bus de dos
hilos simplifica mucho el cableado frente a la alternativa de señales de paso y dirección más
entradas analógicas.

\section{Medida de fuerza con células de carga}

\subsection{Galgas extensométricas y puente de Wheatstone}

Una célula de carga convierte la fuerza en una deformación elástica, que se mide con galgas
extensométricas dispuestas en un puente de Wheatstone \cite{hoffmann2012}. La salida del puente es
una tensión diferencial de unos pocos milivoltios por voltio de excitación a plena escala, que hay
que amplificar con alta ganancia y bajo ruido antes de digitalizarla. Además hay que compensar el
desequilibrio de cero del puente y su deriva con la temperatura \cite{pallas2001}.

\subsection{Acondicionadores integrados}

Existen circuitos integrados que agrupan en un solo chip la excitación, la amplificación, la
conversión A/D y la corrección digital. En este trabajo se han evaluado dos:

\begin{itemize}
  \item \textbf{ZSC31014} de Renesas \cite{zsc31014}. Acondicionador con salida digital
        I\textsuperscript{2}C o SPI, preamplificador de ganancia programable, conversor de 14 bits y
        corrección de ganancia y cero almacenada en EEPROM. Cada lectura incluye dos bits de
        estado que indican si el dato es nuevo o ya se había leído. Su configuración sólo puede
        modificarse en un \emph{modo comando} al que se accede durante una ventana corta tras el
        encendido, lo que ha condicionado mucho el diseño (sección~\ref{sec:prob-zsc}). Se ha usado
        integrado en la placa \emph{Load Cell 2 Click} de MikroElektronika \cite{mikroelc2}.
  \item \textbf{NAU7802} de Nuvoton \cite{nau7802}. Conversor sigma-delta de 24 bits con
        amplificador programable de hasta 128 y regulador interno. Se probó como segunda célula
        (programa \fichero{test\_loadcell\_dual\_v2}) y se retiró del firmware final, que trabaja con
        una sola célula.
\end{itemize}

\section{Microcontrolador}

El proyecto empezó con una placa basada en ESP32-S3 (Seeed XIAO) y se migró a una Raspberry Pi
Pico~2. La tabla~\ref{tab:mcu} resume las características relevantes del RP2350 \cite{rp2350}. La
razón principal del cambio, y sus consecuencias, se explican en la sección~\ref{sec:prob-uart}:
el RP2350 no tiene UART hardware en los pines que usa la placa, pero sí periféricos PIO que pueden
implementar una UART en cualquier pin.

\begin{table}[H]
\centering
\caption{Características del microcontrolador RP2350 en la Raspberry Pi Pico 2 \cite{rp2350}.}
\label{tab:mcu}
\begin{tabular}{@{}ll@{}}
\toprule
Característica & Raspberry Pi Pico 2 (RP2350) \\
\midrule
Núcleos & 2 × Arm Cortex-M33 o 2 × RISC-V Hazard3, a \SI{150}{\mega\hertz} \\
Memoria & \SI{520}{\kilo\byte} SRAM, \SI{4}{\mega\byte} flash QSPI \\
Periféricos serie & 2 × UART, 2 × SPI, 2 × I\textsuperscript{2}C \\
E/S programable & 12 máquinas de estado PIO \\
Entradas analógicas & 3 canales ADC \\
Entorno usado & Arduino-Pico \cite{arduinopico} \\
\bottomrule
\end{tabular}
\end{table}

\section{Buses de comunicación}

\subsection{RS-485 y Modbus RTU}

RS-485 \cite{tia485} es un estándar de capa física diferencial y multipunto, muy extendido en la
industria por su inmunidad al ruido, su alcance y la posibilidad de conectar varios equipos a un
mismo par trenzado. En modo semidúplex, cada nodo activa su emisor sólo mientras transmite; el
control de esa dirección (líneas DE/RE del transceptor) es responsabilidad del maestro.

Modbus \cite{modbusapp} es un protocolo maestro-esclavo de petición y respuesta que modela los
datos del esclavo como registros de 16 bits. En su variante RTU sobre línea serie
\cite{modbusserial}, cada trama lleva la dirección del esclavo, el código de función, los datos y
un CRC-16. Las tramas se delimitan por silencios en la línea de al menos 3,5 caracteres. En este
trabajo se usan la función 0x03 (lectura de registros) y la 0x06 (escritura de un registro), a
través de la biblioteca ModbusMaster \cite{modbusmaster}.

\subsection{I\textsuperscript{2}C}

I\textsuperscript{2}C \cite{i2cspec} es un bus síncrono de dos hilos con resistencias de
\emph{pull-up}, pensado para comunicar circuitos dentro de una misma placa. Es el bus que usan los
acondicionadores de célula de carga. Las placas Click siguen el estándar mikroBUS \cite{mikrobus},
que define un zócalo de 16 pines con I\textsuperscript{2}C, SPI, UART, alimentación y señales
auxiliares.

\section{Supervisión y adquisición en el ordenador}

La supervisión de máquinas se ha hecho tradicionalmente con paneles HMI o paquetes SCADA
comerciales. Para un banco de ensayo de laboratorio, una aplicación web local es una alternativa
ligera: no requiere licencias, funciona en cualquier navegador y permite a varios clientes ver la
máquina a la vez. En este trabajo se ha usado FastAPI \cite{fastapi} sobre el servidor ASGI Uvicorn
\cite{uvicorn}, pySerial \cite{pyserial} para el puerto USB, WebSocket \cite{rfc6455} para el envío
continuo de datos al navegador y Chart.js \cite{chartjs} para las gráficas.

\section{Conclusiones del estado del arte}

Los tambores de radio variable están bien estudiados como principio, y en la literatura se
encuentran métodos de síntesis de su perfil. Lo que motiva este trabajo es disponer de una
plataforma de bajo coste, construida con componentes comerciales, que permita medir en un
prototipo real cuánto aplana el par una caracola y comparar la medida con el modelo. Para ello
hay que integrar un accionamiento industrial por Modbus, un acondicionador de célula de carga con
reglas de configuración estrictas y una cadena de adquisición con marcas de tiempo fiables, que
es lo que desarrollan los capítulos siguientes.
